\scanpage{132}
\mspage{V.3}
% archive sheet number: 131

% L01-04 diagram at top
\diagram{\begin{tikzcd}[ampersand replacement=\&,column sep=small,row sep=large]
\& \& \overset{\deleted{\unsure{\sim}}}{X}'\arrow[dl] \& \overset{\deleted{\sim}}{X}'_{\widetilde\eta}\arrow[l,hook]\arrow[dl] \& \overset{\deleted{\sim}}{U}'\arrow[l,hook]\arrow[dl] \& \overline U'\arrow[l]\arrow[dl] \\
X\arrow[d] \& \widetilde X\arrow[l]\arrow[d] \& \widetilde X_{\widetilde\eta}\arrow[l,hook]\arrow[d] \& \widetilde U\arrow[l,hook]\arrow[d,dashed,no head] \& \overline U\arrow[l]\arrow[d]\arrow[lll,bend left=17,"g"'] \& \\
S \& \widetilde S\arrow[l] \& \widetilde\eta\arrow[l,no head] \& \deleted{\unsure{\overline\eta|\widetilde\eta}}\widetilde\eta\arrow[l] \& \overline\eta\arrow[l] \&
\end{tikzcd}}

% L05-07
un isomorphisme canonique
\[
(\underline\Psi_\nu^{(q)})_{\overline x}\simeq\deleted{H^q(\unsure{\overline U'\times_{\widetilde X}\overline U},\unsure{\Lambda})}\inserted{H^q(\overline U',\Lambda_\nu)},\quad\text{où}\quad\overline U'\simeq X'\times_{\widetilde X}\overline U.
\]

% L08-10
Et comme on a
\[
X'\times_{\widetilde X}\overline U\simeq U'\times_{\widetilde U}\overline U\simeq U'\times_{\widetilde\eta}\overline\eta,\quad-
\]
on trouve finalement
\[
\deleted{\unsure{H^q}}(\underline\Psi_\nu^{(q)})_{\overline x}\simeq H^q(\overline U',\Lambda_\nu),\quad\overline U'=U'\otimes_{\widetilde\eta}\overline\eta=U'\otimes_{\widetilde K}\overline K.
\tag*{\unsure{(1.1.3)}}
\]

% L11-14
Si par exemple $U=X_\eta$ (cas de loin le plus important), on trouve
\[
(\underline\Psi_\nu^{(q)})_{\overline x}\simeq H^q(X'\times_{\widetilde S}\overline\eta,\Lambda_\nu),
\tag*{(1.2)}
\]

% L15-19
i.e. \uline{\unsure{les fibres|la fibre} du faisceau} $\deleted{\unsure{(}}\underline\Psi_\nu^{(q)}\deleted{\unsure{)_{\overline x}}}$ en le point \uline{géom.} $\overline x$ \uline{de} $\widetilde X$ \deleted{sont} \uline{n'est autre que le} \uline{$q$.ième groupe de cohomologie} \inserted{(à coeff. dans $\Lambda_\nu$)} \uline{de la fibre} \inserted{\uline{géométrique}} \uline{générique sur} $\widetilde S$ \uline{du localisé strict} $X'$ \uline{de} $X$ en $\overline x$.

% L20-24 left margin, written slanted along edge
% margin note written slanted
\margin{(De façon imagée, on peut dire que les relations globales entre les deux fibres géométriques $X_{\overline s}$, $X_{\overline\eta}$ se \deleted{\unsure{reflètent|réflètent}}\inserted{\unsure{réfèrent|réduisent}} par les développements précédents aux propriétés des fibres géométriques locales \inserted{\unsure{(i.e. des fibres} génériques du morphisme $X'\longrightarrow\widetilde S$ entre localisés stricts de $X$ et $S$).};)}

% L20 cancelled line above final paragraph
\deleted{\unsure{La même méthode donne} \unsure{pour} les relations,}

% L21-23
Compte tenu des \deleted{\unsure{mêmes}} relations (2.8), on \deleted{\unsure{2}} en conclut des expressions correspondantes pour les

% L24 bottom-edge addition transcribed at the end of the slanted note above;
% the curved leader joins its lower end. Attachment is recorded in the doubts file.

\mspage{V.4}
% archive sheet number: 131A

% L25-26 continuation on second physical sheet
fibres $(\underline\Phi_\nu^{(q)})_{\overline x}$ des faisceaux de cycles évanescents.

% L27-28
\deleted{\unsure{2} Lorsque} Le même argument donne les relations
\[
(\underline\Psi_\nu^{(q)}\deleted{\unsure{)}}_{\mathrm{tame}})_{\overline x}\simeq H^q(U'\otimes_{\widetilde K}\overline K_t),
\tag*{\deleted{\unsure{(2.10)}}}
\]

% L29-31
en particulier si $U=X_\eta$, ceci \deleted{devient} prend la forme
\[
(\underline\Psi_{\nu\,\mathrm{tame}}^q)_{\overline x}=H^q(X\otimes_{\widetilde S}\operatorname{Spec}\overline K_t)
\tag*{\unsure{(1.3)|(1.B)}}
\]
\unsure{;|?} :

% L32 second sheet bottom cancellation
\deleted{\unsure{on a}}
